\chapter{MLC и эквивалентные утверждения}
\noindent В главе формулируется гипотеза MLC и вводятся несколько языков для её количественного изучения: модуль диаметров малых связных множеств, конечные и последовательные аппроксимации, равностепенная аппроксимация и равномерный контроль компонент. Здесь доказываются диагональный критерий и критерий локальной связности через компактные связные множества; связи MLC с равностепенной аппроксимацией и факторизацией устанавливаются в главах 5 и 6.

\section{Обозначения}

Положим
\[
 \N:=\{0,1,2,\ldots\},\qquad f_c(z):=z^2+c,\qquad
 p_0(c):=0,\qquad p_{n+1}(c):=p_n(c)^2+c,
\]
\[
 \M:=\{c\in\C:\sup_{n\in\N}|p_n(c)|<\infty\},\qquad
 K:=\overline B(0,2),\qquad
 \D:=\{x+iy:x,y\in[-2,2]\},
\]
\[
 O_N:=K\cap\bigcap_{j=0}^{N}p_j^{-1}(K),\qquad
 \Pi(z):=\frac{2z}{\max\{2,|z|\}}.
\]
Для непустого \(A\subseteq\C\) и непустых компактных \(A,B\subseteq\C\)
положим
\[
 \dist(z,A):=\inf_{a\in A}|z-a|,\qquad
 \dH(A,B):=
 \max\left\{\sup_{a\in A}\dist(a,B),\
             \sup_{b\in B}\dist(b,A)\right\}.
\]

\section{Основная гипотеза и эквивалентные формулировки}

Пусть \(\mathcal K(\M)\) — множество непустых компактных подмножеств
\(\M\). Для \(C\in\mathcal K(\M)\) положим
\[
 \operatorname{diam}C:=\sup\{|u-v|:u,v\in C\}.
\]
Для \(x,y\in\M\) положим
\[
 \mathcal C_\M(x,y):=
 \{C\in\mathcal K(\M):C\text{ связно},\ x,y\in C\}.
\]
При \(r>0\) положим
\[
 \omega_\M(r):=
  \sup_{\substack{x,y\in\M\\|x-y|<r}}
  \inf_{C\in\mathcal C_\M(x,y)}\operatorname{diam}C
  \in[0,+\infty],\qquad \inf\varnothing:=+\infty.
 \]
 Для \(m\ge1\) положим
 \[
  \mathcal R_m:=
 \left\{G\in C(\D,K):
 \begin{array}{l}
 \exists N\in\N,\ N\ge m,\quad G|_{O_N}=\id_{O_N},\\[1mm]
 \displaystyle\sup_{z\in\D}\dist(G(z),O_N)\le2^{-m}
 \end{array}
 \right\}.
\]
\[
 \mathsf{Fin}:\Leftrightarrow
 \forall m\ge1,\quad\mathcal R_m\ne\varnothing.
 \]
\[
  \begin{gathered}
  \mathsf{EA}_{\mathrm{seq}}\Leftrightarrow
  \exists(m_j)_{j\ge1}\ \exists(G_j)_{j\ge1},\\
  1\le m_1<m_2<\cdots,\qquad m_j\longrightarrow\infty,\qquad
  G_j\in\mathcal R_{m_j}.
  \end{gathered}
 \]
 \[
   \begin{gathered}
   \mathsf{EAdata}\bigl(\rho,(m_j),(G_j)\bigr)\\
  \Leftrightarrow\\
  \begin{gathered}
  1\le m_1<m_2<\cdots,\qquad m_j\longrightarrow\infty,\qquad
  G_j\in\mathcal R_{m_j},\\
 \rho:(0,\infty)\to(0,\infty),\\
  \forall\eta>0,\qquad 0<\rho(\eta)\le\eta,\\
  \forall\eta>0\ \forall j\ge1\ \forall z,w\in\D,\qquad
  |z-w|\le\rho(\eta)\Rightarrow
  |G_j(z)-G_j(w)|\le\eta.
  \end{gathered}
 \end{gathered}
 \]
\[
 \begin{gathered}
 \mathsf{EA}\Leftrightarrow
 \exists\rho\ \exists(m_j)_{j\ge1}\ \exists(G_j)_{j\ge1},\\
 \mathsf{EAdata}\bigl(\rho,(m_j),(G_j)\bigr).
 \end{gathered}
 \]

Для \(F\subseteq\C\) обозначим через \(\pi_0(F)\) множество его
компонент связности. Если \(c\in F\), через
\(\operatorname{Comp}_c(F)\) обозначим компоненту, содержащую \(c\).
При \(c\in\M\), \(r>0\) и \(N\in\N\) положим
\[
 C_N(c,r):=
 \operatorname{Comp}_c\bigl(O_N\cap\overline B(c,r)\bigr).
\]
Для \(L\ge N\) и \(0<\delta<r\) включение
\[
 O_L\cap\overline B(c,\delta)
 \subseteq O_N\cap\overline B(c,r)
\]
индуцирует отображение множеств компонент
\[
 \varphi_{L,N}^{c;\delta,r}:
 \pi_0\bigl(O_L\cap\overline B(c,\delta)\bigr)
 \longrightarrow
 \pi_0\bigl(O_N\cap\overline B(c,r)\bigr),
\]
которое каждую компоненту меньшего множества переводит в единственную
компоненту большего множества, содержащую её.

\begin{definition}[Равномерный контроль компонент]
\label{def:uniform-component-control}
Условие \(\mathsf{UC}\) означает, что
\begin{equation}
 \begin{gathered}
 \forall c\in\M\ \forall\varepsilon>0\
 \exists r,\delta\in\R,\quad 0<\delta<r<\varepsilon,\qquad\\
 \forall N\in\N\ \exists L\in\N,\quad N\le L,\qquad
 \forall A\in\pi_0\bigl(O_L\cap\overline B(c,\delta)\bigr),\\
 \varphi_{L,N}^{c;\delta,r}(A)=C_N(c,r).
 \end{gathered}
 \label{eq:uniform-component-control}
\end{equation}
Радиусы \(r,\delta\) выбираются по \(c,\varepsilon\) и общие для всех
\(N\); индекс \(L\) может зависеть от \(c,\varepsilon,N\).
Для фиксированных \(c,\varepsilon,N,L\) равенство
\(\varphi_{L,N}^{c;\delta,r}(A)=C_N(c,r)\) для всех \(A\) равносильно
включению
\[
 O_L\cap\overline B(c,\delta)\subseteq C_N(c,r).
\]
\end{definition}

 \begin{proposition}[Диагональный критерий выбора с конечным числом оценок]
 \label{prop:diagonal-finite-scale}
 Условие \(\mathsf{EA}\) равносильно следующему условию: существуют
 положительные числа \((\delta_i)_{i\ge1}\), для которых
 \begin{equation}
  \begin{gathered}
  \forall k,M\in\N,\ k,M\ge1\ \exists m\in\N,\ m\ge M\
  \exists G\in\mathcal R_m,\quad\\
  \forall i\in\{1,\ldots,k\}\ \forall z,w\in\D,\qquad
  |z-w|\le\delta_i\Rightarrow
  |G(z)-G(w)|\le2^{-i}.
  \end{gathered}
  \label{eq:diagonal-finite-scale}
 \end{equation}
 \end{proposition}
 \begin{proof}
 Пусть выполнено \(\mathsf{EAdata}(\rho,(m_j),(G_j))\). Положим
 \(\delta_i:=\rho(2^{-i})\). Для любых \(k,M\ge1\) выберем \(j\ge k\)
 так, чтобы \(m_j\ge M\). Тогда \(G_j\in\mathcal R_{m_j}\), а оценка
 равностепенной непрерывности при \(\eta=2^{-i}\) даёт все неравенства
 в \eqref{eq:diagonal-finite-scale} для \(1\le i\le k\).

 Обратно, пусть существуют \((\delta_i)\), удовлетворяющие
 \eqref{eq:diagonal-finite-scale}. Положим \(m_0:=0\). Последовательно
 применяя это условие с \(k=j\) и \(M=m_{j-1}+1\), выберем
 \(m_j\ge m_{j-1}+1\) и \(G_j\in\mathcal R_{m_j}\), для которых
 \[
  \forall i\in\N,\ 1\le i\le j\ \forall z,w\in\D,\qquad
  |z-w|\le\delta_i\Rightarrow
  |G_j(z)-G_j(w)|\le2^{-i}.
 \]
 Тогда \(m_1<m_2<\cdots\), и \(m_j\to\infty\). Зафиксируем
 \(\eta>0\) и выберем \(i\ge1\) так, чтобы \(2^{-i}\le\eta\).
 Равномерная непрерывность конечного семейства
 \(G_1,\ldots,G_{i-1}\) даёт \(\sigma_\eta>0\), для которого
 \[
  |z-w|\le\sigma_\eta\Rightarrow
  |G_j(z)-G_j(w)|\le\eta
  \qquad(1\le j<i,\ z,w\in\D);
 \]
 при \(i=1\) положим \(\sigma_\eta:=1\). Положим
 \(\rho(\eta):=\min\{\eta,\delta_i,\sigma_\eta\}\).
 Если \(j\ge i\), требуемая оценка следует из неравенства с индексом
 \(i\); если \(j<i\), она следует из выбора \(\sigma_\eta\).
 Следовательно, \(\mathsf{EAdata}(\rho,(m_j),(G_j))\), то есть
 \(\mathsf{EA}\).
 \end{proof}
 \begin{equation}
  \forall\rho,(m_j),(G_j),\qquad
   \mathsf{EAdata}\bigl(\rho,(m_j),(G_j)\bigr)
 \Rightarrow
  \forall\eta>0,\qquad
  \omega_\M\!\left(\frac{\rho(\eta)}2\right)\le2\eta.
  \label{eq:quantitative-diameter-bound}
\end{equation}
\begin{equation}
  \begin{gathered}
  \mathsf{Fin},\qquad \mathsf{EA}_{\mathrm{seq}},\\
  \mathsf{EA}\quad\Leftrightarrow\quad
  \lim_{r\downarrow0}\omega_\M(r)=0
  \quad\Leftrightarrow\quad\LC(\M).
  \end{gathered}
 \label{eq:main-discrepancy}
\end{equation}

Теорема об аппроксимации из главы~\ref{ch:approximation}
даёт требуемую в \(\mathsf{EA}_{\mathrm{seq}}\) последовательность:
\(m_j=j\), \(G_j=T_j\in\mathcal R_j\). Это отдельный результат о конечных
приближениях; отображения \(T_j\) не входят в доказательство
\(\mathsf{AM}\), \(\mathsf{GC}\) или условной факторизации. Для редукций используется лишь диадическая сетка, построенная в
главе~\ref{ch:approximation}.
Условие \(\mathsf{EA}\) усиливает существование таких отображений
требованием общего модуля равностепенной непрерывности \(\rho\).
По предложению~\ref{prop:diagonal-finite-scale} диагональный путь
требует для каждого конечного \(k\) выполнения первых \(k\) оценок
осцилляции при сколь угодно больших \(m\). Теорема
\ref{thm:uniform-component-control} устанавливает
\(\mathsf{UC}\Longleftrightarrow\LC(\M)\), откуда следует, что
безусловная факторизация локальных компонент уже влечёт MLC. Обратная
импликация от локальной связности к \(\mathsf{AM}\) или \(\mathsf{GC}\)
здесь не рассматривается и не требуется для целевой цепочки. Ниже
условная теорема доказывает \(\mathsf H_{\M}\Rightarrow\mathsf{UC}\), а
оставшиеся строгие редукции разделены по двум открытым оценкам.
Искомая логическая цепочка имеет вид
\[
 \mathsf H_{\M}\;(\mathsf{AM}+\mathsf{GC})
 \ \Longrightarrow\ \mathsf{UC}
 \ \Longleftrightarrow\ \LC(\M).
\]
Итак, доказательство \(\mathsf{AM}\) и \(\mathsf{GC}\) завершило бы
доказательство MLC через теорему о факторизации; обратные импликации к
открытым оценкам здесь не предполагаются.

\section{Критерий через компактные связные множества малого диаметра}

\begin{theorem}[Критерий локальной связности через компактные связные множества]
\label{thm:small-connected-compacts}
Пусть \(X\) — компактное метрическое пространство с метрикой \(d\).
Тогда \(X\) локально связно тогда и только тогда, когда
\[
 \forall\varepsilon>0\ \exists\delta>0\ \forall x,y\in X,\qquad
 d(x,y)<\delta\ \Rightarrow\
 \exists C\subseteq X:
 \begin{cases}
 C\text{ компактно и связно},\\
 x,y\in C,\\
 \operatorname{diam}(C)<\varepsilon.
 \end{cases}
\]
\end{theorem}
\begin{proof}
Если \(X\) локально связно, то у каждой точки есть связная открытая
окрестность диаметра меньше \(\varepsilon\). Компактность даёт
конечное подпокрытие этих окрестностей, а лемма Лебега — число
\(\delta>0\). Если \(d(x,y)<\delta\), точки \(x,y\) лежат в одной
такой окрестности; её замыкание в \(X\) компактно, связно и имеет
диаметр меньше \(\varepsilon\).

Обратно, пусть \(x\in V\), где \(V\) открыто в \(X\). Выберем
\(\varepsilon>0\) так, чтобы \(B(x,\varepsilon)\cap X\subseteq V\).
Для соответствующего \(\delta\) каждая точка \(y\in B(x,\delta)\cap X\)
лежит вместе с \(x\) в связном множестве \(C\) диаметра меньше
\(\varepsilon\), поэтому \(C\subseteq V\). Значит, компонента \(V\),
содержащая \(x\), содержит шар \(B(x,\delta)\cap X\). Компоненты
открытых множеств открыты, что равносильно локальной связности.
\end{proof}

\begin{proposition}[Критерий локальной связности через функцию \(\omega_\M\)]
\label{prop:omega-criterion}
\[
 \lim_{r\downarrow0}\omega_\M(r)=0
 \quad\Leftrightarrow\quad
 \LC(\M).
\]
\end{proposition}
\begin{proof}
Пусть \(\lim_{r\downarrow0}\omega_\M(r)=0\). Для каждого
\(\varepsilon>0\) существует \(\delta>0\) такое, что
\(\omega_\M(\delta)<\varepsilon\). При \(|x-y|<\delta\) имеем
\[
 \inf_{C\in\mathcal C_\M(x,y)}\operatorname{diam}C
 \le\omega_\M(\delta)<\varepsilon.
\]
Значит, существует \(C\in\mathcal C_\M(x,y)\) с
\(\operatorname{diam}C<\varepsilon\). Теорема
\ref{thm:small-connected-compacts} даёт \(\LC(\M)\).

Обратно, пусть \(\LC(\M)\). Применим теорему
\ref{thm:small-connected-compacts} к \(\varepsilon/2\) и получим
\(\delta>0\). При \(0<r\le\delta\) для каждой пары
\(x,y\in\M\), \(|x-y|<r\), существует
\(C\in\mathcal C_\M(x,y)\) с \(\operatorname{diam}C<\varepsilon/2\).
Следовательно, \(\omega_\M(r)\le\varepsilon/2<\varepsilon\), откуда
\(\lim_{r\downarrow0}\omega_\M(r)=0\).
\end{proof}
